% Generated draft: requires mathematical and source review.
\documentclass[11pt]{article}
\usepackage[margin=25mm]{geometry}
\usepackage{fontspec}
\setmainfont{Latin Modern Roman}
\usepackage{amsmath,amssymb,mathtools,unicode-math}
\setmathfont{Latin Modern Math}
\usepackage{enumitem,needspace,xurl,hyperref}
\setlength{\parindent}{0pt}
\setlength{\parskip}{5pt}
\setlength{\emergencystretch}{3em}
\newcommand{\N}{\mathbb{N}}
\newcommand{\Z}{\mathbb{Z}}
\newcommand{\Q}{\mathbb{Q}}
\newcommand{\R}{\mathbb{R}}
\newcommand{\C}{\mathbb{C}}
\newcommand{\F}{\mathbb{F}}
\newcommand{\A}{\mathbb{A}}
\newcommand{\PP}{\mathbb{P}}
\newcommand{\E}{\mathbb{E}}
\newcommand{\eps}{\varepsilon}
\newcommand{\dd}{\,\mathrm{d}}
\begin{document}
\section*{UnsolvedMath ID: 6200089}
\textbf{Source problem code:} catalogue entry
\subsection{Definitions and mathematical statement}
Let $\Delta$ be a developable triangle of groups. Developability means that its local groups and inclusions are realized by stabilizers of an action of its complex-of-groups fundamental group $\Gamma=\pi_1(\Delta)$ on a simply connected development, with quotient the triangle. Thus $\pi_1(\Delta)$ is not the ordinary fundamental group of the underlying contractible triangle.

A discrete group has Kazhdan's Property $(T)$ if every unitary representation with almost invariant unit vectors has a nonzero invariant vector. Almost invariant means that, for every finite subset $F$ of the group and $\varepsilon>0$, some unit vector $v$ satisfies $\|\rho(g)v-v\|<\varepsilon$ for all $g\in F$.

Suppose every cell group of $\Delta$ has Property $(T)$, and every link graph $L$ in the development has
\[
\lambda_1(L)>\frac12.
\]
Does it follow that $\Gamma$ has Property $(T)$? This preserves both the cell-group hypothesis and the strict spectral threshold.

The spectral normalization is missing. In a usual finite connected unweighted link convention, $\lambda_1$ is the first positive eigenvalue of the normalized graph Laplacian
\[
(\mathcal Lf)(v)=f(v)-\frac1{\deg(v)}\sum_{w\sim v}f(w),
\]
considered with degree-weighted inner product. Weighted or infinite links need corresponding operator and spectral-gap definitions. The source gives no local-finiteness, connectivity, weighting, or normalization hypotheses, so the displayed finite-link convention is context, not an added source assumption. Its phrase ``universal cover of T'' has no separately defined space $T$; the development of $\Delta$ is the apparent intended reference and needs checking.
\subsection{Short English statement}
Must a developable triangle of Property $(T)$ groups have Property $(T)$ fundamental group when every development link has spectral gap greater than $1/2$? The link operator and development conventions are unspecified.
\subsection{Sources}
\begin{itemize}
\item Local UnsolvedMath registry, permanent ID 6200089, exact source code catalogue entry, ``Boundaries of Groups and Kleinian Groups---Problem 89''; proposed year 2005, no proposer supplied. Registry status open, classification OPEN-TRIAGE.
\item Kapovich, \emph{Problems on Boundaries of Groups and Kleinian Groups} (2005), Problem 89, PDF page 22, \url{https://www.math.ucdavis.edu/~kapovich/EPR/problems.pdf}. Supplied normalized pdf-text extraction accessed July 29, 2026. No PDF was accessed here.
\item The supplied summary relates the question to complexes of groups and Żuk-type spectral criteria, but supplies no precise theorem or publication and reports the triangle-of-groups assertion unestablished. That is a supplied open assessment, not an independent current-status check.
\item Provenance flags current status and statement-reuse terms NEEDS\_REVIEW; the document is publicly author-hosted.
\end{itemize}
\end{document}
